\chapter{Giới thiệu}

Chương này trình bày bối cảnh và động lực nghiên cứu của khóa luận, khảo sát các
phương pháp kiểm thử bảo mật ứng dụng web hiện có, đánh giá các công cụ tiêu biểu,
phân tích đặc thù của lỗ hổng Blind SSRF, xác định khoảng trống còn tồn tại và đề
xuất hướng giải quyết.
Phần~\ref{sec:motivation} mô tả bối cảnh và lý do chọn đề tài.
Phần~\ref{sec:methods-overview} khảo sát các phương pháp kiểm thử bảo mật chính.
Phần~\ref{sec:existing-tools} đánh giá các công cụ tiêu biểu hiện có.
Phần~\ref{sec:ssrf-challenge} phân tích đặc thù của lỗ hổng SSRF và Blind SSRF.
Phần~\ref{sec:gap} xác định khoảng trống nghiên cứu và hướng tiếp cận đề xuất.
Phần~\ref{sec:objectives} phát biểu mục tiêu và phạm vi.
Phần~\ref{sec:contributions} liệt kê các đóng góp kỹ thuật của khóa luận.
Phần~\ref{sec:structure} trình bày cấu trúc khóa luận.

% ============================================================
\section{Bối cảnh và Lý do Chọn Đề tài}
\label{sec:motivation}
% ============================================================

\subsection{Tầm quan trọng của Ứng dụng Web và Bảo mật}

Ứng dụng web ngày nay là hạ tầng số không thể thiếu của hầu hết các tổ chức. Từ
ngân hàng trực tuyến, thương mại điện tử, đến hệ thống quản lý nhà nước và hồ sơ y
tế điện tử, logic nghiệp vụ cốt lõi đều được hiện thực hóa dưới dạng web application
chạy trên nền HTTP/HTTPS. Sự phức tạp ngày càng tăng của các hệ thống này kéo theo
bề mặt tấn công rộng hơn: một nền tảng SaaS hiện đại có thể tích hợp hàng chục
microservice nội bộ, gọi tới nhiều API bên thứ ba, xử lý tệp đính kèm từ người dùng,
và tương tác với cơ sở dữ liệu với đặc quyền cao. Mỗi điểm tích hợp đó đều là một
bề mặt tấn công tiềm năng.

PHP là ngôn ngữ lập trình phổ biến nhất phía server trong hệ sinh thái web, được sử
dụng bởi hơn 77\% các website có ngôn ngữ server-side đã biết tính đến năm
2023~\cite{w3techs2023php}, bao gồm các nền tảng lớn như WordPress (chiếm hơn
43\% website toàn cầu), Drupal, Magento và Joomla. Hệ sinh thái plugin của WordPress
--- với hơn 60.000 plugin từ bên thứ ba, mỗi plugin là một bề mặt tấn công tiềm năng
--- là nguồn sinh ra liên tục các lỗ hổng bảo mật mới được công bố dưới dạng CVE
(Common Vulnerabilities and Exposures). Sự kết hợp giữa mức độ phổ biến cao của PHP
và hệ sinh thái plugin thiếu kiểm soát tạo ra môi trường màu mỡ cho kẻ tấn công,
đồng thời đặt ra thách thức lớn cho kiểm thử bảo mật tự động.

Trong bối cảnh đó, kiểm thử bảo mật tự động --- đặc biệt là kỹ thuật \textit{fuzzing}
--- đã trở thành một hướng nghiên cứu quan trọng, được áp dụng rộng rãi trong cả môi
trường học thuật lẫn công nghiệp~\cite{manes2021art, zhu2022fuzzing}. Tuy nhiên, như
phân tích trong các phần tiếp theo sẽ chỉ ra, một lớp lỗ hổng nguy hiểm và ngày càng
phổ biến --- Blind SSRF --- vẫn nằm ngoài tầm quan sát của toàn bộ thế hệ công cụ
kiểm thử tự động hiện tại.

\subsection{SSRF và Vị trí trong OWASP Top 10}

Server-Side Request Forgery (SSRF) là lớp lỗ hổng bảo mật trong đó kẻ tấn công
khiến máy chủ ứng dụng thực hiện các yêu cầu mạng tùy ý tới địa chỉ do kẻ tấn công
kiểm soát~\cite{owasp2021top10}. SSRF đã leo lên vị trí thứ mười trong danh sách
OWASP Top 10:2021 --- lần đầu tiên được đưa vào danh sách này với tư cách một mục
riêng biệt, không ghép với các lỗ hổng khác. OWASP ghi nhận rằng dù tỷ lệ xuất hiện
của SSRF còn khiêm tốn so với SQLi hay XSS, tỷ lệ khai thác thành công lại cao hơn
đáng kể, vì SSRF thường khai thác tin tưởng ngầm giữa các thành phần trong hạ tầng
nội bộ --- máy chủ web thường được phép truy cập các dịch vụ nội bộ như Redis,
Elasticsearch, hay cloud metadata endpoint mà không qua xác thực, vì chúng vốn không
được thiết kế để tiếp xúc trực tiếp với Internet.

Tác động của SSRF đặc biệt nghiêm trọng trong kiến trúc đám mây. Sự cố Capital One
năm 2019 --- một trong những vụ rò rỉ dữ liệu ngân hàng lớn nhất lịch sử với hơn
100 triệu hồ sơ khách hàng bị ảnh hưởng --- được kích hoạt thông qua SSRF tới
instance metadata service (IMDS) của AWS (địa chỉ \texttt{169.254.169.254}), cho
phép kẻ tấn công lấy IAM credentials tạm thời và truy cập S3 bucket chứa dữ liệu
nhạy cảm~\cite{owasp2023ssrf}. Sự cố này cho thấy SSRF không chỉ là lỗ hổng lý
thuyết mà có thể dẫn đến hậu quả thực tế nghiêm trọng ở quy mô lớn.

\subsection{Blind SSRF --- Điểm mù của Kiểm thử Tự động}

Dạng nguy hiểm và khó phát hiện nhất của SSRF là \textit{Blind SSRF}: máy chủ thực
hiện network request nhưng không phản ánh response về cho người dùng --- ứng dụng
xử lý kết quả nội bộ, loại bỏ, hoặc chỉ ghi log. Từ góc độ giao thức HTTP, mọi thứ
có vẻ bình thường: không có error message, không có status code bất thường, không có
nội dung khác biệt trong response. Dấu hiệu duy nhất tồn tại ở tầng mạng --- máy
chủ đã khởi tạo một kết nối ra ngoài tới địa chỉ mà kẻ tấn công kiểm soát. Điều
này đặc biệt phổ biến trong các kịch bản: xác nhận webhook (ứng dụng gọi URL để
kiểm tra tính hợp lệ nhưng không trả về nội dung), background processing (SSRF được
kích hoạt bởi job bất đồng bộ), hay tích hợp dịch vụ (ứng dụng tự động lấy metadata
từ URL người dùng cung cấp).

Chính đặc điểm vô hình này khiến Blind SSRF trở thành điểm mù lớn nhất của hầu hết
công cụ kiểm thử bảo mật tự động hiện có. Không một greybox fuzzer nào ---
Phuzz~\cite{neef2024phuzz}, webFuzz~\cite{vanrooij2021webfuzz}, hay
Witcher~\cite{trickel2023witcher} --- có cơ chế để quan sát hành vi mạng của ứng
dụng mục tiêu vượt ra ngoài HTTP response và code coverage. Khóa luận này xuất phát
từ quan sát đó và đề xuất một hướng tiếp cận có hệ thống để lấp khoảng trống: tích
hợp kênh tín hiệu \textit{Out-of-Band} (OOB) vào vòng lặp greybox fuzzing, biến một
máy chủ lắng nghe bên ngoài thành một kênh phản hồi độc lập cho fuzzer.

\begin{figure}[htbp]
    \centering
    % Mô tả hình: Sơ đồ luồng tấn công SSRF điển hình.
    % Bên trái: Attacker gửi HTTP request chứa URL độc hại (url=http://169.254.169.254/...).
    % Ở giữa: Web Server / Ứng dụng PHP nhận request, gọi file_get_contents($url).
    % Bên phải: Các dịch vụ nội bộ bị tấn công: AWS Metadata (169.254.169.254),
    %           Redis (:6379), Elasticsearch (:9200), Admin panel (/admin).
    % Mũi tên: Attacker → Web App (HTTP/HTTPS public), Web App → Internal (qua tường lửa nội bộ).
    % Ghi chú: Firewall cho phép web app kết nối nội bộ nhưng chặn truy cập trực tiếp từ Internet.
    \includegraphics[width=0.88\textwidth]{figures/fig-ssrf-flow}
    \caption{Luồng tấn công SSRF.}
    \label{fig:ssrf-flow}
\end{figure}

% ============================================================
\section{Phương pháp Kiểm thử Bảo mật Ứng dụng Web}
\label{sec:methods-overview}
% ============================================================

Kiểm thử bảo mật ứng dụng web đã phát triển qua nhiều thế hệ công cụ và phương
pháp~\cite{li2018fuzzing, manes2021art}. Mỗi cách tiếp cận có điểm mạnh và hạn chế
riêng, đặc biệt khi đối mặt với các lớp lỗ hổng không để lại dấu vết hiển thị như
Blind SSRF.

\subsection{Phân tích Tĩnh}

Phân tích tĩnh (static analysis) kiểm tra mã nguồn mà không cần thực thi chương
trình. Đối với PHP, phân tích tĩnh có thể truy vết luồng dữ liệu (taint analysis)
từ điểm nhập (source) như \texttt{\$\_GET['url']} tới điểm sink nguy hiểm như
\texttt{file\_get\_contents(\$url)}. Phân tích tĩnh là bước đầu tiên hữu ích trong
pipeline kiểm thử bảo mật: chi phí thấp, không yêu cầu triển khai ứng dụng, và có
thể tích hợp vào quy trình CI/CD.

Tuy nhiên, phân tích tĩnh gặp khó khăn với \textit{context-sensitive taint tracking}
qua nhiều tầng framework: trong một WordPress plugin, đầu vào người dùng thường đi
qua nhiều lớp abstraction trước khi đến sink. Phân tích tĩnh không thể đánh giá hành
vi runtime --- một URL có bị lọc hay không phụ thuộc vào logic điều kiện và cấu hình
runtime, không thể suy luận từ code tĩnh. Kết quả là tỷ lệ false positive cao và
false negative đáng kể với các lỗ hổng phức tạp.

\subsection{Thực thi Tượng trưng và Concolic Execution}

Thực thi tượng trưng (symbolic execution) thay thế các giá trị cụ thể bằng biểu thức
tượng trưng và giải các ràng buộc để tìm đầu vào kích hoạt mọi nhánh
code~\cite{godefroid2008automated}. Kỹ thuật này có khả năng đạt độ phủ lý thuyết
cao và xác định chính xác điều kiện kích hoạt lỗ hổng. Concolic execution kết hợp
thực thi cụ thể và tượng trưng~\cite{ganesh2009taint}, giảm bớt vấn đề ``path
explosion'' nhưng vẫn tốn kém về mặt tính toán.

Với ứng dụng web PHP thực tế --- hàng nghìn dòng code, framework phức tạp, tương tác
database, và nhiều dependency bên thứ ba --- chi phí tính toán của symbolic execution
trở nên không khả thi. Ngoài ra, các lời gọi hàm ngoài (như \texttt{curl\_exec})
phụ thuộc vào môi trường runtime thực tế (phản hồi từ OOB server), điều mà symbolic
execution không thể mô phỏng mà không có môi trường thực thi cụ thể.

\subsection{Kiểm thử Hộp đen (Blackbox)}

Kiểm thử blackbox không có bất kỳ thông tin nào về cấu trúc bên trong ứng dụng ---
công cụ chỉ tương tác qua giao thức HTTP, quan sát đầu ra để phát hiện bất thường.
Ưu điểm là dễ triển khai, không cần mã nguồn, và áp dụng được cho bất kỳ ứng dụng
web nào. Nhược điểm cơ bản là khả năng khám phá code path bị giới hạn bởi những gì
quan sát được từ response. Doupe và cộng sự~\cite{doupe2010johnny} đã chứng minh
rằng blackbox scanner bỏ sót đáng kể các lỗ hổng đòi hỏi nhiều bước tương tác hoặc
phụ thuộc vào trạng thái session phức tạp.

\subsection{Kiểm thử Greybox và Coverage-Guided Fuzzing}

Greybox fuzzing khai thác thông tin thực thi thu thập được trong quá trình chạy ---
chủ yếu là code coverage --- làm tín hiệu phản hồi để định hướng quá trình sinh đột
biến~\cite{fioraldi2020afl, zalewski2014afl}. AFL (American Fuzzy Lop) và sau đó là
AFL++~\cite{fioraldi2020afl} đã chứng minh tính hiệu quả của cách tiếp cận này cho
ứng dụng nhị phân. Các công trình gần đây nhất về kiểm thử bảo mật web tự động ---
webFuzz~\cite{vanrooij2021webfuzz}, Witcher~\cite{trickel2023witcher}, và
Phuzz~\cite{neef2024phuzz} --- đều theo hướng greybox fuzzing vì đây là cân bằng tốt
nhất giữa chi phí và hiệu quả: không yêu cầu mã nguồn đầy đủ như whitebox, không bị
giới hạn bởi oracle response như blackbox, và mở rộng tốt hơn symbolic execution.

Bảng~\ref{tab:method-comparison} so sánh bốn phương pháp theo các tiêu chí quan
trọng với bài toán phát hiện Blind SSRF.

\begin{table}[h]{So sánh các phương pháp kiểm thử bảo mật theo các tiêu chí liên quan
đến phát hiện Blind SSRF}
\label{tab:method-comparison}
\centering
\begin{tabular}{|l|c|c|c|c|}
\hline
\textbf{Tiêu chí} & \textbf{Tĩnh} & \textbf{Symbolic} & \textbf{Blackbox} & \textbf{Greybox} \\
\hline
Không cần mã nguồn     & \ding{55} & \ding{55} & \checkmark & \checkmark \\
Khả năng mở rộng       & Tốt & Kém & Tốt & Tốt \\
Phủ code path sâu      & Hạn chế & Cao & Thấp & Trung bình–cao \\
Phát hiện runtime sink & \ding{55} & Có thể & \checkmark & \checkmark \\
Tích hợp kênh OOB      & \ding{55} & \ding{55} & Bán thủ công & \textbf{Có thể} \\
\hline
\end{tabular}
\end{table}

Bảng~\ref{tab:method-comparison} cho thấy greybox fuzzing là nền tảng phù hợp nhất
để tích hợp kênh OOB: nó đã có vòng lặp tự động, có khả năng thu thập tín hiệu
runtime, và không phụ thuộc vào mã nguồn. Cần thiết kế thêm kênh OOB song song với
coverage --- đây là đóng góp kỹ thuật của khóa luận này.

% ============================================================
\section{Thực trạng Công cụ Kiểm thử Bảo mật Hiện có}
\label{sec:existing-tools}
% ============================================================

\subsection{Công cụ Blackbox}

\textbf{OWASP ZAP (Zed Attack Proxy)~\cite{zapdevteam2023}} là công cụ kiểm thử
bảo mật web mã nguồn mở được sử dụng rộng rãi nhất, cung cấp cả chế độ thủ công và
tự động hóa qua API và daemon mode. ZAP thực hiện spider để khám phá endpoint, sau đó
gửi payload chèn sẵn và phân tích HTTP response để phát hiện SQLi, XSS, path
traversal, và một số lớp lỗ hổng khác. ZAP không tích hợp kênh OOB, do đó không thể
phát hiện Blind SSRF.

\textbf{Burp Suite Professional~\cite{portswigger2023burp}} cung cấp tính năng Burp
Collaborator --- một infrastructure OOB cho phép phát hiện Blind SSRF, Blind SQLi, và
Blind XXE trong kiểm thử bán thủ công. Khi kiểm thử viên chèn Collaborator payload
vào tham số, bất kỳ interaction nào từ ứng dụng mục tiêu tới Collaborator server đều
được ghi nhận và báo cáo. Đây là tính năng mạnh nhưng đòi hỏi sự chủ động của kiểm
thử viên --- chọn endpoint, đặt payload, theo dõi kết quả --- không phải tự động hóa
hoàn toàn theo nghĩa của fuzzing.

\subsection{Công cụ Greybox Fuzzing cho Web}

\textbf{webFuzz~\cite{vanrooij2021webfuzz}} (ESORICS 2021) là một trong những greybox
fuzzer web đầu tiên được công bố. webFuzz thu thập coverage bằng cách instrument
JavaScript phía client và HTML response, tập trung chủ yếu vào phát hiện XSS.
webFuzz không hỗ trợ phát hiện các lỗ hổng server-side như SSRF, SQLi hay command
injection do không có instrumentation phía server.

\textbf{Witcher~\cite{trickel2023witcher}} (IEEE S\&P 2023) mở rộng greybox fuzzing
cho PHP bằng cách sử dụng taint analysis để theo dõi luồng dữ liệu từ input tới các
sink nguy hiểm. Witcher đạt kết quả đáng kể cho SQLi và command injection nhưng hoàn
toàn không có cơ chế phát hiện SSRF vì không có hook trên các hàm mạng và không có
kênh OOB.

\textbf{Phuzz~\cite{neef2024phuzz}} (AsiaCCS 2024) là greybox fuzzer PHP toàn diện
nhất hiện tại, hỗ trợ phát hiện sáu loại lỗ hổng server-side (SQLi, RCE, path
traversal, insecure deserialization, XXE, open redirect) và XSS phía client. Trong
bộ benchmark 87 lỗ hổng, Phuzz đạt tỷ lệ phát hiện 99\%. Tuy nhiên, như chính bài
báo gốc thừa nhận, SSRF --- đặc biệt Blind SSRF --- hoàn toàn nằm ngoài khả năng
phát hiện của Phuzz. Bài báo liệt kê SSRF là mục tiêu ưu tiên cao nhất trong phần
Future Work, với ghi chú cụ thể: \textit{``The SSRF could have been found by Phuzz,
if \texttt{wp\_remote\_get} was hooked and included in the vulnerability detection.''}

\subsection{Tổng hợp So sánh}

Bảng~\ref{tab:tool-comparison} tổng hợp khả năng phát hiện SSRF và Blind SSRF của
các công cụ tiêu biểu, làm rõ khoảng trống mà khóa luận này hướng tới lấp đầy.

\begin{table}[h]{So sánh khả năng phát hiện Blind SSRF của các công cụ kiểm thử
bảo mật ứng dụng web tiêu biểu}
\label{tab:tool-comparison}
\centering
\begin{tabular}{|l|l|c|c|c|}
\hline
\textbf{Công cụ} & \textbf{Loại} & \textbf{SSRF} & \textbf{Blind SSRF} & \textbf{Tự động} \\
\hline
OWASP ZAP~\cite{zapdevteam2023}       & Blackbox & Hạn chế   & \ding{55}   & \checkmark \\
Burp Suite~\cite{portswigger2023burp} & Blackbox & \checkmark & \checkmark & Bán tự động \\
webFuzz~\cite{vanrooij2021webfuzz}    & Greybox  & \ding{55}  & \ding{55}  & \checkmark \\
Witcher~\cite{trickel2023witcher}      & Greybox  & \ding{55}  & \ding{55}  & \checkmark \\
Phuzz~\cite{neef2024phuzz}            & Greybox  & \ding{55}  & \ding{55}  & \checkmark \\
\textbf{Phuzz+OOB (đề xuất)}          & Greybox  & \checkmark & \checkmark & \checkmark \\
\hline
\end{tabular}
\end{table}

Bảng~\ref{tab:tool-comparison} cho thấy rõ khoảng trống: không có greybox fuzzer nào
có thể phát hiện Blind SSRF một cách tự động. Burp Suite có khả năng này nhưng đòi
hỏi can thiệp thủ công đáng kể; tất cả greybox fuzzer hoàn toàn mù với SSRF vì thiếu
kênh tín hiệu OOB.

% ============================================================
\section{Lỗ hổng SSRF và Thách thức Phát hiện Tự động}
\label{sec:ssrf-challenge}
% ============================================================

\subsection{Nguyên lý và Phân loại SSRF}

Lỗ hổng SSRF phát sinh khi ứng dụng web nhận URL hay hostname từ đầu vào người dùng
và dùng nó làm tham số cho một hàm mạng phía server mà không xác thực đầy đủ. Kẻ
tấn công có thể truyền URL trỏ tới:

\begin{itemize}
    \item \textbf{Hệ thống nội bộ}: \texttt{http://192.168.1.1/admin} --- tường lửa
    chặn truy cập từ Internet nhưng không chặn từ chính máy chủ ứng dụng, vì máy chủ
    ứng dụng được coi là tin cậy trong mạng nội bộ.
    \item \textbf{Cloud metadata endpoint}: \texttt{http://169.254.169.254/latest/\\meta-data/iam/security-credentials/}
    --- đọc IAM credentials của AWS EC2 instance, dẫn tới chiếm quyền toàn bộ tài
    khoản cloud.
    \item \textbf{Dịch vụ nội bộ không xác thực}: Redis, Elasticsearch, Jenkins,
    Kibana thường không yêu cầu xác thực khi truy cập từ mạng nội bộ, vì chúng không
    được thiết kế để tiếp xúc trực tiếp với Internet.
    \item \textbf{Leo thang thành RCE}: Kết hợp SSRF với Gopher protocol và Redis RESP
    injection để ghi SSH key hoặc thực thi lệnh shell tùy ý trên máy chủ.
\end{itemize}

Theo mức độ quan sát được từ bên ngoài, SSRF chia thành hai dạng chính:
\textit{Regular SSRF} (in-band) khi response của request phía server được phản ánh
một phần hoặc toàn bộ trong HTTP response trả về cho người dùng; và \textit{Blind
SSRF} khi ứng dụng thực hiện request nhưng không trả kết quả ra ngoài --- xử lý nội
bộ, ghi log, hoặc chỉ kiểm tra trạng thái kết nối.

\subsection{Kỹ thuật Bypass Filter SSRF}

Nhiều ứng dụng cố gắng ngăn chặn SSRF bằng blocklist địa chỉ IP nội bộ
(127.0.0.1, 10.0.0.0/8, 192.168.0.0/16, 172.16.0.0/12). Kẻ tấn công có thể vượt
qua bằng nhiều kỹ thuật:

\begin{itemize}
    \item \textbf{Biểu diễn IP thay thế}: \texttt{0x7f000001} (hex),
    \texttt{2130706433} (decimal), \texttt{0177.0.0.1} (octal) đều biểu diễn
    127.0.0.1 nhưng thường qua được kiểm tra regex đơn giản.
    \item \textbf{IPv6}: \texttt{::1} hay \texttt{[::ffff:127.0.0.1]} thường bị bỏ
    sót trong blocklist chỉ kiểm tra IPv4.
    \item \textbf{URL encoding}: \texttt{\%31\%32\%37.\%30.\%30.\%31} hay
    double-encoding có thể bypass kiểm tra chuỗi thô.
    \item \textbf{DNS-based redirect}: Domain của attacker trỏ về IP nội bộ; hoặc
    dùng các dịch vụ như \texttt{nip.io}/\texttt{sslip.io} --- ví dụ
    \texttt{127.0.0.1.nip.io} phân giải thành \texttt{127.0.0.1}.
    \item \textbf{DNS rebinding}: Domain ban đầu trỏ về IP public để qua kiểm tra;
    sau khi TTL hết, DNS rebind thành địa chỉ nội bộ. Kỹ thuật này bypass cả các
    kiểm tra IP được thực hiện trước khi gọi hàm mạng.
    \item \textbf{Userinfo trick}: \texttt{http://trusted.com@172.20.0.10:8000/} ---
    một số parser nhầm \texttt{172.20.0.10} là đường dẫn, không phải host.
    \item \textbf{Scheme-less}: Ứng dụng tự prepend \texttt{https://} vào giá trị
    tham số; bypass bằng cách chỉ cung cấp host không có scheme.
    \item \textbf{Fragment bypass}: \texttt{http://172.20.0.10\#trusted.com} ---
    fragment bị bỏ qua khi gửi request thực nhưng có thể đánh lừa validator.
\end{itemize}

Sự đa dạng của các kỹ thuật bypass là lý do tại sao một bộ mutator SSRF hiệu quả cần
bao phủ nhiều biến thể, không chỉ plain IP. Chi tiết cụ thể về cài đặt các biến thể
này trong \texttt{SSRFMutator} được trình bày ở Chương~3.

\subsection{Tại sao Tín hiệu Truyền thống Không đủ}

Bảng~\ref{tab:signal-comparison} so sánh các loại tín hiệu phản hồi mà công cụ kiểm
thử tự động hiện tại có thể khai thác và khả năng áp dụng cho từng loại lỗ hổng.

\begin{table}[h]{Tín hiệu phản hồi và khả năng phát hiện từng loại lỗ hổng
(\checkmark = có tín hiệu, \ding{55} = không có tín hiệu)}
\label{tab:signal-comparison}
\centering
\begin{tabular}{|l|c|c|c|c|}
\hline
\textbf{Tín hiệu phản hồi} & \textbf{SQLi} & \textbf{XSS} & \textbf{SSRF} & \textbf{Blind SSRF} \\
\hline
Nội dung HTTP response          & \checkmark & \checkmark & \checkmark & \ding{55} \\
HTTP status code bất thường     & \checkmark & \ding{55}  & Có thể     & \ding{55} \\
Exception / thông báo lỗi       & \checkmark & \ding{55}  & Có thể     & \ding{55} \\
Độ phủ mã (code coverage)       & \checkmark & \checkmark & \checkmark & \ding{55} \\
\textbf{OOB network callback}   & \ding{55}  & \ding{55}  & \checkmark & \checkmark \\
\hline
\end{tabular}
\end{table}

\begin{figure}[htbp]
    \centering
    % Mô tả hình: So sánh hai kịch bản SSRF cạnh nhau (hai cột).
    % Cột trái — Regular SSRF: Attacker → Web App → Internal Service → response nội dung
    %   được phản ánh về Attacker trong HTTP response. Mũi tên khép kín.
    % Cột phải — Blind SSRF: Attacker → Web App → Internal Service → response được xử lý
    %   nội bộ hoặc ghi log; Attacker nhận HTTP 200 OK bình thường, không có nội dung bất thường.
    %   Chỉ OOB Server (bên dưới, tách biệt) ghi nhận được kết nối từ Web App.
    % Điểm nhấn: với Blind SSRF, tất cả tín hiệu truyền thống (response, status, coverage)
    %   đều bình thường; chỉ OOB callback mới là bằng chứng.
    \includegraphics[width=0.92\textwidth]{figures/fig-blind-vs-regular}
    \caption{So sánh Regular SSRF và Blind SSRF.}
    \label{fig:blind-vs-regular}
\end{figure}

Hình~\ref{fig:blind-vs-regular} và Bảng~\ref{tab:signal-comparison} cùng làm rõ vấn
đề cốt lõi: với Blind SSRF, mọi tín
hiệu truyền thống đều vắng mặt. Không có nội dung bất thường trong response, không
có status code lạ, không có exception. Đặc biệt, code path thực thi hoàn toàn giống
nhau dù URL trỏ tới OOB server hay một URL bình thường --- \texttt{file\_get\_contents}
hay \texttt{curl\_exec} đều được gọi đúng một lần, không có nhánh code mới nào được
kích hoạt. OOB network callback là tín hiệu \emph{duy nhất} có khả năng xác nhận
Blind SSRF, và đây là lý do tại sao không có greybox fuzzer nào hiện tại có thể phát
hiện được loại lỗ hổng này.

\subsection{Đặc thù PHP và Bề mặt Tấn công SSRF}

PHP cung cấp một bề mặt tấn công SSRF đặc biệt rộng do nhiều hàm đa năng có thể
thực hiện HTTP request ngầm khi nhận URL làm đối số:

\begin{itemize}
    \item \textbf{Sink rõ ràng}: \texttt{curl\_exec()}, \texttt{fsockopen()},
    \texttt{stream\_socket\_client()} --- hàm mạng tường minh, thường được kiểm thử
    viên chú ý.
    \item \textbf{Sink đa năng}: \texttt{file\_get\_contents()}, \texttt{fopen()},
    \texttt{readfile()} --- vốn dùng để đọc file cục bộ nhưng âm thầm thực hiện
    HTTP/FTP request khi nhận URL.
    \item \textbf{Sink không trực quan}: \texttt{getimagesize()} --- thường dùng để
    lấy kích thước ảnh, ít ai ngờ đây là SSRF sink. Thực tế lỗ hổng
    CVE-2022-1751 (Skitter Slideshow) khai thác đúng sink này.
    \item \textbf{Sink WordPress}: \texttt{wp\_remote\_get()},
    \texttt{wp\_remote\_post()} --- tầng abstraction cao hơn, được nhiều plugin dùng
    thay cho cURL, tạo ra sink gián tiếp không được bao phủ bởi hook tầng thấp nếu
    không instrument riêng.
\end{itemize}

Ngoài ra, kịch bản Command Injection dẫn tới SSRF cũng đáng chú ý: khi input người
dùng được truyền vào hàm shell như \texttt{shell\_exec()}, kẻ tấn công có thể chèn
lệnh \texttt{curl} hoặc \texttt{wget} để thực hiện kết nối mạng tùy ý. Đây là một
dạng SSRF gián tiếp qua tầng shell, cần bộ mutator và hook riêng biệt.

% ============================================================
\section{Khoảng trống Nghiên cứu và Hướng Giải quyết}
\label{sec:gap}
% ============================================================

\subsection{Xác định Khoảng trống}

Phân tích ở các phần trước xác định một khoảng trống cụ thể và rõ ràng trong công
nghệ kiểm thử bảo mật hiện tại: \textit{không tồn tại greybox fuzzer nào tích hợp
kênh tín hiệu OOB vào vòng lặp fuzzing để phát hiện Blind SSRF một cách tự động và
không có false positive}.

Khoảng trống này tồn tại do ba nguyên nhân có liên quan:

\begin{enumerate}
    \item \textbf{Giả định thiết kế}: Greybox fuzzer truyền thống được thiết kế với
    giả định rằng tín hiệu phản hồi đến từ hai kênh: HTTP response và code coverage.
    Cả hai kênh đều không phản ánh hành vi mạng ẩn của ứng dụng.
    \item \textbf{Khoảng cách công nghệ}: Kỹ thuật OOB tồn tại trong công cụ kiểm
    thử thủ công (Burp Collaborator) nhưng chưa được tích hợp vào vòng lặp fuzzing
    tự động theo cách có hệ thống, với cơ chế tương quan token đảm bảo chính xác
    trong môi trường song song.
    \item \textbf{Thách thức kỹ thuật}: Phát hiện Blind SSRF trong fuzzing tự động
    đặt ra bài toán không tầm thường --- cần tương quan nhân quả giữa hàng nghìn
    request gửi song song và các OOB callback không đồng bộ, đồng thời tránh false
    positive từ các kết nối mạng không liên quan.
\end{enumerate}

\subsection{Hướng Giải quyết Đề xuất}

Khóa luận này đề xuất lấp khoảng trống đó bằng cách mở rộng Phuzz~\cite{neef2024phuzz}
với ba thành phần mới được tích hợp theo cách không phá vỡ kiến trúc gốc:

\begin{enumerate}
    \item \textbf{Bộ đột biến SSRF chuyên biệt.} \texttt{SSRFMutator} sinh payload
    OOB bao phủ 21 biến thể (IP-based, DNS-based, hostname-based) theo phương pháp
    round-robin, đảm bảo đa dạng kỹ thuật bypass. \texttt{CmdInjectionSSRFMutator}
    thêm 10 biến thể command injection cho các endpoint có nguy cơ gọi lệnh shell.
    Thiết kế đa biến thể là cần thiết vì một số lỗ hổng chỉ phát hiện được qua đúng
    một biến thể: ví dụ CVE-2020-28975 (Canto plugin) chỉ phát hiện qua biến thể
    scheme-less, CVE-2021-32682 (elFinder) chỉ qua biến thể hostname-based.

    \item \textbf{Máy chủ OOB lắng nghe độc lập.} Container \texttt{phuzz-oob} chạy
    Flask HTTP/HTTPS listener (cổng 8000/8443) và dnsmasq với wildcard DNS
    (\texttt{*.oob}), nhận callback từ ứng dụng mục tiêu và ghi log nguyên tử vào
    shared-tmpfs. Thiết kế container độc lập cho phép tái sử dụng với bất kỳ ứng dụng
    mục tiêu nào mà không cần sửa đổi.

    \item \textbf{Module phát hiện hai tầng.} \texttt{SSRFVulnCheck} kết hợp tín hiệu
    in-band (PHP hook qua UOPZ ghi nhận lời gọi hàm mạng với OOB payload) và tín hiệu
    OOB (callback log từ OOB server) theo logic AND. Chỉ khi cả hai file bằng chứng
    cùng tồn tại mới xác nhận lỗ hổng, đảm bảo zero false positive.
\end{enumerate}

Hướng tiếp cận này giải quyết đồng thời cả ba nguyên nhân của khoảng trống: bổ sung
kênh tín hiệu thứ ba (OOB), tự động hóa hoàn toàn quá trình từ sinh payload đến xác
nhận lỗ hổng, và đảm bảo tương quan chính xác qua cơ chế token 8 ký tự gắn với từng
ứng viên.

% ============================================================
\section{Mục tiêu và Phạm vi Nghiên cứu}
\label{sec:objectives}
% ============================================================

\subsection{Mục tiêu}

Mục tiêu chính của khóa luận là thiết kế và cài đặt một hệ thống phát hiện Blind
SSRF hoàn toàn tự động trong môi trường greybox fuzzing PHP, tích hợp vào framework
Phuzz, với các tiêu chí đánh giá cụ thể:

\begin{itemize}
    \item \textbf{Tự động hóa hoàn toàn}: không yêu cầu can thiệp thủ công trong quá
    trình fuzzing; hệ thống tự sinh payload, thu tín hiệu từ cả hai kênh, tương quan
    OOB callback với request, xác nhận lỗ hổng và ghi report có cấu trúc.
    \item \textbf{Không có false positive}: mỗi lỗ hổng chỉ được báo cáo khi có bằng
    chứng mạng thực tế từ cả hai tầng tín hiệu (logic AND).
    \item \textbf{Tương thích và không xâm lấn}: tích hợp hoàn toàn dưới dạng module
    bổ sung, giữ nguyên toàn bộ chức năng phát hiện lỗ hổng gốc và cấu trúc code của
    Phuzz.
    \item \textbf{Bao phủ bề mặt tấn công}: payload bao gồm 21 biến thể SSRF và 10
    biến thể command injection SSRF; PHP hook bao phủ 23 hàm trên 5 nhóm chức năng.
    \item \textbf{Kiểm chứng thực nghiệm}: đánh giá trên cả ứng dụng lab và CVE thực
    tế từ hệ sinh thái WordPress plugin.
\end{itemize}

\subsection{Phạm vi}

\textbf{Trong phạm vi:}
\begin{itemize}
    \item Ứng dụng PHP (phiên bản 7.4 và 8.x) chạy trong môi trường Docker, bao gồm
    cả ứng dụng WordPress với plugin bên thứ ba.
    \item Lỗ hổng SSRF và Blind SSRF được kích hoạt qua tham số HTTP (query string,
    POST form data, JSON body, HTTP header tùy chỉnh).
    \item Phát hiện SSRF thông qua các hàm mạng PHP: \texttt{file\_get\_contents},
    \texttt{curl\_exec}, \texttt{fsockopen}, \texttt{getimagesize}, các hàm WordPress
    HTTP API, và các hàm shell execution (qua command injection chèn lệnh curl/wget).
    \item Môi trường kiểm thử cục bộ (Docker network 172.20.0.0/24); OOB server
    trong cùng mạng nội bộ với địa chỉ IP tĩnh 172.20.0.10.
\end{itemize}

\textbf{Ngoài phạm vi:}
\begin{itemize}
    \item Các ngôn ngữ server-side khác (Python, Node.js, Java, Ruby). Khóa luận
    thảo luận khả năng mở rộng nhưng không cài đặt.
    \item Stored SSRF --- trong đó payload được lưu vào database và kích hoạt SSRF ở
    một request hoàn toàn khác sau đó. Khóa luận phân tích thách thức và đề xuất
    hướng giải quyết ở Chương~5 nhưng không cài đặt đầy đủ.
    \item Khai thác lỗ hổng (exploitation) sau khi phát hiện; khóa luận chỉ tập
    trung vào giai đoạn phát hiện và xác nhận với bằng chứng mạng.
    \item Phát hiện phân tán trên nhiều máy vật lý hoặc môi trường cloud production
    với tường lửa egress nghiêm ngặt.
\end{itemize}

% ============================================================
\section{Đóng góp của Khóa luận}
\label{sec:contributions}
% ============================================================

Khóa luận có ba đóng góp kỹ thuật chính:

\begin{enumerate}
    \item \textbf{Bộ đột biến SSRF chuyên biệt.} \texttt{SSRFMutator} sinh
    payload OOB bao phủ 21 biến thể (IP-based, DNS-based, hostname-based)
    theo phương pháp round-robin, bao phủ các kỹ thuật bypass phổ biến nhất.
    \texttt{CmdInjectionSSRFMutator} bổ sung 10 biến thể command injection
    cho các endpoint có nguy cơ thực thi lệnh shell. Đây là bộ đột biến SSRF
    đầu tiên được tích hợp vào vòng lặp greybox fuzzing tự động cho ứng
    dụng PHP.

    \item \textbf{Cơ sở hạ tầng OOB lắng nghe tích hợp.} Container
    \texttt{phuzz-oob} cung cấp HTTP/HTTPS listener (cổng 8000/8443) và
    dnsmasq với wildcard DNS (\texttt{*.oob}), nhận callback từ ứng dụng
    mục tiêu và ghi log nguyên tử vào shared-tmpfs. Thiết kế container độc
    lập cho phép tái sử dụng với bất kỳ ứng dụng mục tiêu nào mà không cần
    sửa đổi.

    \item \textbf{Cơ chế phát hiện Blind SSRF hai tầng không có false positive.}
    \texttt{SSRFVulnCheck} kết hợp tín hiệu in-band (PHP hook qua UOPZ ghi
    nhận lời gọi hàm mạng với OOB payload) và tín hiệu OOB (callback log
    từ OOB server) theo logic AND. Chỉ khi cả hai bằng chứng cùng tồn tại
    mới xác nhận lỗ hổng, đảm bảo zero false positive ngay trong môi trường
    fuzzing song song.
\end{enumerate}

Ba đóng góp trên được tích hợp hoàn toàn theo mô hình plugin không xâm lấn
vào Phuzz~\cite{neef2024phuzz}, giữ nguyên toàn bộ chức năng phát hiện lỗ
hổng gốc, đồng thời mở rộng khả năng phát hiện lên lớp lỗ hổng Blind SSRF
vốn nằm ngoài tầm quan sát của mọi greybox fuzzer hiện tại.

% ============================================================
\section{Cấu trúc Khóa luận}
\label{sec:structure}
% ============================================================

Phần còn lại của khóa luận được tổ chức như sau:

\textbf{Chương~2 --- Cơ sở Lý thuyết} trình bày nền tảng lý thuyết cần thiết: tổng
quan chi tiết về kỹ thuật greybox fuzzing, energy-based seed scheduling và coverage
guidance; kiến trúc và cơ chế hoạt động đầy đủ của framework Phuzz bao gồm cơ chế
instrumentation và VulnChecker; cơ chế phân tích động, hook hàm PHP qua extension
UOPZ; và nhận xét về bài toán cần giải quyết từ góc độ khoảng trống kỹ thuật.

\textbf{Chương~3 --- Thiết kế và Cài đặt Hệ thống} mở đầu bằng phân tích toàn diện
lỗ hổng SSRF (định nghĩa, phân loại, tác động, kỹ thuật bypass) và kỹ thuật OOB
(cơ chế tương quan token, so sánh giao thức), tiếp theo là mô tả chi tiết kiến trúc
bốn container Docker trên mạng phuzz-net 172.20.0.0/24; cài đặt hook PHP với 23 hook
trên 5 nhóm hàm; OOB server (Flask + dnsmasq, ghi log nguyên tử); bộ đột biến SSRF
với 21 biến thể và 10 biến thể command injection; cơ chế phát hiện hai tầng với AND
logic; và cách các thành phần tích hợp vào vòng lặp fuzz chính theo mô hình plugin
không xâm lấn.

\textbf{Chương~4 --- Thực nghiệm và Đánh giá} trình bày môi trường thực nghiệm, thiết
kế thực nghiệm trên 14 mục tiêu (6 ứng dụng lab và 8 CVE thực tế từ hệ sinh thái
WordPress và PHP plugin), kết quả phát hiện chi tiết cho từng mục tiêu, phân tích các
trường hợp đặc biệt (biến thể bypass scheme-less, sink getimagesize, stream wrapper
cho include/require), và so sánh với Phuzz gốc chưa tích hợp OOB.

\textbf{Chương~5 --- Kết luận và Hướng Phát triển} tổng kết các đóng góp chính của
khóa luận, phân tích trung thực các hạn chế còn tồn tại (phạm vi ngôn ngữ PHP, stored
SSRF, quy mô thực nghiệm, phụ thuộc kết nối mạng nội bộ), và đề xuất các hướng mở
rộng tiềm năng (đa ngôn ngữ, stored SSRF qua deferred lookup, XXE/SQLi OOB, tự động
tạo PoC, tích hợp CI/CD pipeline).
